\section{Introduction}
Large language models (LLMs) excel at reasoning, coding, and multimodal tasks. However, due to the high hardware requirements (RAM, GPU, etc.) they demand, these models are usually run only on cloud-based infrastructures or very advanced local systems. This creates a significant bottleneck for system developers seeking to reduce cloud dependency, ensure data privacy, and avoid high operational costs.

As a result of these limitations, small language models (SLMs) are emerging as a strong alternative. Their ability to run on local machines provides a significant advantage in terms of data security, and their operational costs are much lower compared to LLMs. The ability to fine-tune them for specific domains, their easy integration with explicit control modules, and their collaborative operation structure make them quite attractive. However, SLMs cannot provide the broad expertise offered by a single massive model (like GPT-4o). They still fall short in tasks requiring logical reasoning, deep contextual understanding, and extensive world knowledge. Therefore, reducing the issue to a simple "SLM or LLM?" competition is no longer a fair or comprehensive approach. Instead, current research revolves around questions such as in which areas a single SLM is sufficient, when an LLM is needed, when a hybrid or agent-based architecture should be used, and how a customized SLM system can outperform a single LLM.

The existing literature provides concrete examples of how these architectural choices are shaped by the dynamics of the field. For example, in the telecommunications sector, the Symbiotic Agents study has managed to significantly reduce error rates and resource consumption by using hybrid and multi-agent systems \cite{chatzistefanidis2025symbiotic}. In medicine, the Self-Reflective Chest X-Ray Report Generation study combined three different AI models in a single pipeline---a self-reflective agent (LLaMA-8B) and an LLM judge---to create a hybrid structure that maximized diagnostic accuracy and minimized error rates \cite{han2025self}. Similarly, in cybersecurity, Pietrzak and colleagues solved a complex fake report detection problem using only a fine-tuned single SLM (TinyLlama) and detection models (TinyBERT, DistilBERT) \cite{pietrzak2025what}. They achieved this success with just standard consumer-grade GPU hardware, achieving higher detection rates compared to large models like GPT-4 \cite{pietrzak2025what}.

These developments provide practical alternatives for scenarios where relying on a single large language model is insufficient or excessively costly. Moreover, research in this area not only focuses on basic accuracy rates but also aims to provide tangible improvements in computational load, operational cost, and latency. Researchers have clearly demonstrated which types of hybrid designs work for specific tasks and what kinds of architectures are effective in real-world conditions.

This review investigates how hybrid models converge across different fields, when LLMs should play a more active role, and in which cases control can be delegated to SLMs. It examines in which specific domains fine-tuned SLMs perform best and where they naturally fall short. By analyzing recurring hybrid methods in the literature, the article reveals scenarios where SLMs can realistically replace or complement large language models. Since the scope is fundamentally built on the resource-performance trade-off, this review places special emphasis on hybrid and multi-agent systems.

The main contributions of this review are as follows:
\begin{itemize}
    \item It provides a synthesized overview of 54 studies covering cloud-based monolithic LLMs, hybrid SLM-LLM architectures, the performance of fine-tuned SLMs versus large models, and multi-agent workflows.
    \item It frames the literature around four distinct practical research questions: Performance, Efficiency, Architecture, and Domain Suitability.
    \item It creates an analytical review matrix that reveals the commonalities and differences among the studies, mapping which specific research question each study addresses.
\end{itemize}

\section{Findings}
\subsection{Corpus-Level Trends Before the RQ Analysis}
Before answering the four research questions directly, three corpus-level observations are worth emphasizing.

First, the field is now primarily a systems literature rather than a pure model-scaling literature. Although a few studies still compare standalone small and large models directly, the dominant trend is to embed them inside a larger control strategy. Hybrid edge-cloud pipelines, retrieval-centered systems, expert routers, task decomposers, self-reflective workflows, and fusion modules are now common design choices. In other words, recent work increasingly treats model scale as one variable in a broader systems problem.

Second, hybridization is motivated by different constraints in different domains. In telecom and networking, the main concerns are deployment, context fidelity, and latency budgets. In healthcare and scientific tasks, the dominant concerns are reliability, interpretability, and domain specificity. In industrial monitoring, multimodal hybridization is driven by the need to combine coarse global understanding with fine-grained local detail. In optimisation and education, hybrid methods often serve as a way to recover reasoning quality from small models through feedback, routing, or targeted prompting.

Third, the evidence base is uneven. Accuracy is almost always reported; compute, latency, and cost are not. This matters because a large share of the literature frames hybridization as an efficiency solution, yet only a minority of studies expose sufficiently comparable numbers to test those claims rigorously. The result is a literature that often proves that hybridization \emph{can} work, but less often proves how much it costs in practice under standardized conditions.

\subsection{RQ1: Comparative Performance of Hybrid and Multi-Agent Systems}
Across the corpus, hybrid and multi-agent architectures generally improve the performance ceiling of SLM-centered systems and, in several cases, also outperform larger monolithic baselines. The main reason is not that small models suddenly become universally strong, but that coordination mechanisms route the right subtasks, tokens, or contexts to the right component at the right moment.

The clearest evidence comes from edge-cloud collaboration. Hao et al.~\cite{hao2024hybrid} show that token-level draft-verification can recover LLM-like quality on GSM8K while using only 25.8\% of the full-cloud LLM cost. Their abstract and introduction clarify why this works: many tokens are easy enough for the on-device SLM, and only a small subset of difficult tokens must be verified or regenerated in the cloud. The key performance implication is that quality recovery does not require whole-query escalation. Token-level selectivity is enough.

This pattern also appears in wireless and mobile settings. Oh et al.~\cite{oh2025uncertainty} frame hybrid inference as a distributed speculative-decoding problem and show that uncertainty-aware skipping can preserve up to 97.54\% of LLM accuracy. Their work is especially important for the present review because it ties performance directly to a measurable systems lever: if the SLM can estimate when its token is likely to be accepted, then the system can skip unnecessary communication and cloud inference without losing much quality.

Routing can do more than merely preserve LLM quality. In dialogue state tracking, OrchestraLLM uses a retrieval-based router to select the most suitable expert among SLM and LLM pools~\cite{lee2024orchestrallm}. The website extraction reports 82.46\% turn-level joint goal accuracy on MultiWOZ and 68.09\% on SGD, outperforming both the LLM and SLM baselines. The corresponding PDF makes the reason more explicit: SLMs and LLMs exhibit complementary strengths, with SLMs handling schema and formatting constraints well and LLMs contributing commonsense and broader contextual reasoning. This is an example of a hybrid system exceeding both ends of the scale spectrum because the system exploits complementarity rather than averaging behavior.

Domain-specific retrieval produces similarly strong results. TeleOracle adapts Phi-2 to telecom question answering by combining fine-tuning, semantic chunking, hybrid retrieval, reranking, and context-window extension~\cite{alabbasi2024teleoracle}. In the reviewed corpus, this system reaches 81.20\% accuracy, which is roughly 30\% above the base Phi-2 model and competitive with much